
%%
%% This is file `sample-sigplan.tex',
%% generated with the docstrip utility.
%%
%% The original source files were:
%%
%% samples.dtx  (with options: `all,proceedings,bibtex,sigplan')
%% 
%% IMPORTANT NOTICE:
%% 
%% For the copyright see the source file.
%% 
%% Any modified versions of this file must be renamed
%% with new filenames distinct from sample-sigplan.tex.
%% 
%% For distribution of the original source see the terms
%% for copying and modification in the file samples.dtx.
%% 
%% This generated file may be distributed as long as the
%% original source files, as listed above, are part of the
%% same distribution. (The sources need not necessarily be
%% in the same archive or directory.)
%%
%%
%% Commands for TeXCount
%TC:macro \cite [option:text,text]
%TC:macro \citep [option:text,text]
%TC:macro \citet [option:text,text]
%TC:envir table 0 1
%TC:envir table* 0 1
%TC:envir tabular [ignore] word
%TC:envir displaymath 0 word
%TC:envir math 0 word
%TC:envir comment 0 0
%%
%% The first command in your LaTeX source must be the \documentclass
%% command.
%%
%% For submission and review of your manuscript please change the
%% command to \documentclass[manuscript, screen, review]{acmart}.
%%
%% When submitting camera ready or to TAPS, please change the command
%% to \documentclass[sigconf]{acmart} or whichever template is required
%% for your publication.
%%
%%
\documentclass[sigplan,screen]{acmart}
%%
%% \BibTeX command to typeset BibTeX logo in the docs
\AtBeginDocument{%
  \providecommand\BibTeX{{%
    Bib\TeX}}}

%% Rights management information.  This information is sent to you
%% when you complete the rights form.  These commands have SAMPLE
%% values in them; it is your responsibility as an author to replace
%% the commands and values with those provided to you when you
%% complete the rights form.

\setcopyright{acmlicensed}
\copyrightyear{2026}
\acmYear{2026}
\acmDOI{XXXXXXX.XXXXXXX}
%% These commands are for a PROCEEDINGS abstract or paper.
\acmConference[CSCI 5922 Spring 2026]{Fundamentals of Neural Networks and Deep Learning}{April 2026}{Boulder, CO}
%%
%%  Uncomment \acmBooktitle if the title of the proceedings is different
%%  from ``Proceedings of ...''!
%%
%%\acmBooktitle{Woodstock '18: ACM Symposium on Neural Gaze Detection,
%%  June 03--05, 2018, Woodstock, NY}
\acmISBN{978-1-4503-XXXX-X/2026/02}


%%
%% Submission ID.
%% Use this when submitting an article to a sponsored event. You'll
%% receive a unique submission ID from the organizers
%% of the event, and this ID should be used as the parameter to this command.
%%\acmSubmissionID{123-A56-BU3}

%%
%% For managing citations, it is recommended to use bibliography
%% files in BibTeX format.
%%
%% You can then either use BibTeX with the ACM-Reference-Format style,
%% or BibLaTeX with the acmnumeric or acmauthoryear sytles, that include
%% support for advanced citation of software artefact from the
%% biblatex-software package, also separately available on CTAN.
%%
%% Look at the sample-*-biblatex.tex files for templates showcasing
%% the biblatex styles.
%%

%%
%% The majority of ACM publications use numbered citations and
%% references.  The command \citestyle{authoryear} switches to the
%% "author year" style.
%%
%% If you are preparing content for an event
%% sponsored by ACM SIGGRAPH, you must use the "author year" style of
%% citations and references.
%% Uncommenting
%% the next command will enable that style.
%%\citestyle{acmauthoryear}

\usepackage{listings}
\usepackage{xcolor}
\usepackage{graphicx}
\usepackage{placeins}
\usepackage{algorithm}
\usepackage{algpseudocode}
\lstset{
  basicstyle=\ttfamily\small,
  breaklines=true,
  breakatwhitespace=true,
  postbreak=\mbox{\textcolor{red}{$\hookrightarrow$}\space},
  frame=single,
  columns=fullflexible
}

\usepackage{subcaption} % add this in your preamble
\usepackage{tikz}
\usetikzlibrary{arrows.meta,positioning,backgrounds}
\usetikzlibrary{shapes.geometric,fit,calc}

\usepackage{forest}

%%
%% end of the preamble, start of the body of the document source.
\begin{document}

%%
%% The "title" command has an optional parameter,
%% allowing the author to define a "short title" to be used in page headers.
\title[Interpreting Personality Prediction]{Interpreting Personality Prediction: Word-Level vs Sentence-Level Contributions in Neural Models}

%%
%% The "author" command and its associated commands are used to define
%% the authors and their affiliations.
%% Of note is the shared affiliation of the first two authors, and the
%% "authornote" and "authornotemark" commands
%% used to denote shared contribution to the research.

\author{Anna Mathew}
\email{Anna.Mathew@colorado.edu}
\affiliation{%
  \country{}
  \institution{University of Colorado Boulder}
}
\author{Rutuja Nikumb}
\email{RutujaRavindra.Nikumb@colorado.edu}
\affiliation{%
  \country{}
  \institution{University of Colorado Boulder}
}


%%
%% By default, the full list of authors will be used in the page
%% headers. Often, this list is too long, and will overlap
%% other information printed in the page headers. This command allows
%% the author to define a more concise list
%% of authors' names for this purpose.
\renewcommand{\shortauthors}{Anna Mathew and Rutuja Nikumb}

% Remove header/footer on page 1. Running heads continue from page 2.
\fancypagestyle{firstpagestyle}{%
  \fancyhf{}
  \renewcommand{\headrulewidth}{0pt}
  \renewcommand{\footrulewidth}{0pt}
}

%%
%% The abstract is a short summary of the work to be presented in the
%% article.

\begin{abstract}
This project studies interpretability reliability in neural personality prediction on MBTI-500 text data. We compare attention-, gradient-, and perturbation-based explanations at word and sentence levels to test agreement across methods, local versus global encoding behavior, and sensitivity to spurious topic/style cues. The outcome is a reproducible evaluation plan that moves beyond classification accuracy and tests whether explanation claims are stable and meaningful.
\end{abstract}

\begin{CCSXML}
<ccs2012>
 <concept>
  <concept_id>10003120.10003121.10003129</concept_id>
  <concept_desc>Agentic Approach</concept_desc>
  <concept_significance>300</concept_significance>
 </concept>
</ccs2012>
\end{CCSXML}

\keywords{personality prediction, interpretability, MBTI, NLP, deep learning}

\received{2026}

%%
%% This command processes the author and affiliation and title
%% information and builds the first part of the formatted document.


\maketitle

% Custom running heads: no conference/location text in headers.
\pagestyle{fancy}
\fancyhf{}
\fancyhead[LE]{\small\itshape\shortauthors}
\fancyhead[RO]{\small\itshape\shorttitle}
\fancyfoot[C]{\thepage}

\section{Introduction}
Accurate personality prediction from text can support adaptive learning systems, mental-health triage support, personalized recommendation, and human-computer interaction design. However, in high-stakes settings, predictive performance alone is not enough; we also need interpretable evidence for why a model predicts a personality label. Our project is motivated by this gap between model accuracy and explanation quality for personality inference from user-generated text.

Existing personality prediction systems largely optimize classification metrics while offering limited analysis of what textual evidence drives predictions. Even when interpretability methods are applied, studies often report one explanation method in isolation and mostly at a single granularity (token-level or feature-level), making it difficult to assess reliability. Prior work also does not sufficiently test whether models use meaningful personality cues versus spurious correlations such as topic words or social-media style markers.

We propose a comparative interpretability study on the MBTI-500 dataset \cite{mbti500_kaggle} using transformer-based neural models \cite{devlin2019bert}. Our central research question is: \textit{How do neural models attribute importance to words versus sentences when predicting personality, and are these attributions consistent across explanation methods?} We will compare attention-, gradient-, and perturbation-based explanations, test local (word) versus global (sentence) encoding hypotheses, and evaluate robustness to spurious cues. This is new in combining cross-method agreement with multi-granularity ablations and spurious-signal tests in one unified experimental pipeline.

\section{Related Work}
\subsection{Topic 1: Neural Personality Prediction from Text}
Majumder et al. \cite{majumder2017deep} introduce deep document modeling for personality detection and show neural models can outperform traditional feature-engineering baselines on text-based personality tasks. Gjurkovi\'c and \v{S}najder \cite{gjurkovic2018reddit} demonstrate that Reddit language provides useful behavioral signals for personality prediction and highlight the value of social-media corpora for this problem. Majumder et al. \cite{majumder2020deep} further emphasize end-to-end deep learning for personality inference, reinforcing that learned representations improve predictive quality over shallow approaches.

Our work is different because it focuses on \textit{interpretability reliability} (agreement, granularity sensitivity, and spurious-cue dependence) rather than only predictive performance.

\subsection{Topic 2: Feature Attribution and Explanation Methods in NLP}
Ribeiro et al. \cite{ribeiro2016lime} propose LIME, a perturbation-based local explanation method that approximates model behavior around an instance to identify influential features. Sundararajan et al. \cite{sundararajan2017axiomatic} introduce Integrated Gradients, a gradient-based attribution method grounded in axioms such as sensitivity and implementation invariance. Lundberg and Lee \cite{lundberg2017unified} provide SHAP, a unified attribution framework based on Shapley values with strong consistency guarantees and broad model applicability.

Our work is different because we apply these explanation families jointly to personality prediction and explicitly quantify \textit{cross-method consistency} for both words and sentences.

\subsection{Topic 3: Reliability of Attention and Faithfulness of Explanations}
Jain and Wallace \cite{jain2019attention} show that attention distributions are often weakly correlated with other feature-importance measures and can be manipulated without major prediction changes. Serrano and Smith \cite{serrano2019attention} test attention faithfulness via erasure analyses and report that high attention weights do not always correspond to causally important tokens, challenging the assumption that attention is intrinsically explanatory.

Our work is different because we do not evaluate one explanation channel in isolation; we directly test where methods agree or disagree and connect disagreement to downstream behavior via removal experiments.

\subsection{Topic 4: Spurious Correlations and Shortcut Learning}
Gururangan et al. \cite{gururangan2018annotation} identify annotation artifacts that allow models to exploit superficial cues, demonstrating how dataset biases can inflate apparent benchmark performance. Geirhos et al. \cite{geirhos2020shortcut} formalize shortcut learning and show that deep models often rely on easy but non-robust patterns instead of task-relevant semantics. Sagawa et al. \cite{sagawa2020investigation} analyze spurious correlations under overparameterization and show that standard empirical risk minimization can fail on minority or distribution-shifted groups.

Our work is different because we operationalize shortcut analysis for personality prediction by removing topic-heavy and stylistic markers, then measuring prediction shifts and performance changes.

\section{Methods}
\subsection{Problem Setup and Data}
We use the MBTI-500 dataset \cite{mbti500_kaggle}, where each instance contains text and one of 16 MBTI labels. We frame this as a 16-class classification task and also report binary splits per trait dimension (E/I, S/N, T/F, J/P) for additional analysis.

\subsection{Model and Training Configuration}
We fine-tune a BERT-base encoder \cite{devlin2019bert} with a classification head. Text is segmented into sentences, tokenized with WordPiece, and truncated/padded to a fixed maximum length. We perform train/validation/test splits with fixed random seeds, use early stopping on validation macro-F1, and report mean and standard deviation across 3 runs. Hyperparameters (learning rate, batch size, max length, epochs) are selected by validation performance.

\subsection{Explanation Pipeline}
For each test sample, we compute:
\begin{itemize}
    \item \textbf{Attention-based importance:} aggregated attention scores for tokens and sentence-level pooled scores.
    \item \textbf{Gradient-based importance:} token-level integrated gradients \cite{sundararajan2017axiomatic}, then sentence scores by averaging token attributions within each sentence.
    \item \textbf{Perturbation-based importance:} prediction change after masking top candidate words or sentences.
\end{itemize}

Importance scores are normalized per sample, then ranked to obtain top-$k$ words and top-$k$ sentences. This enables method-to-method comparison at both granularity levels.

\subsection{Research Questions}
\begin{itemize}
    \item \textbf{RQ1:} Do different interpretability methods (attention, gradients, perturbation) agree on what parts of text are important for personality prediction?
    \item \textbf{RQ2:} Is personality encoded in localized signals (specific words) or distributed across higher-level structures (sentences)?
    \item \textbf{RQ3:} Are models relying on genuine personality cues or spurious correlations such as topics or slang?
    \item \textbf{RQ4:} Do human judgments of important words/sentences align with model explanations?
\end{itemize}

\section{Experiments}
\subsection{Experiment 1: Cross-Method Agreement Analysis}
\textbf{Main purpose:} Quantify whether attention, gradient, and perturbation explanations are consistent about which words/sentences matter for personality prediction. This validates whether interpretability conclusions are method-dependent.
This experiment directly evaluates \textbf{RQ1} by measuring agreement across interpretability methods.

\textbf{Procedure:}
\begin{itemize}
    \item For each method and sample, extract top-$k$ words and top-$k$ sentences.
    \item Compute pairwise agreement between methods at word and sentence levels.
    \item Aggregate by class and overall dataset.
\end{itemize}

\textbf{Evaluation metrics:}
\begin{itemize}
    \item \textbf{Jaccard similarity} on top-$k$ sets (overlap of highlighted evidence).
    \item \textbf{Spearman correlation} on full ranked importance scores (global ranking agreement).
\end{itemize}
Jaccard captures discrete overlap, while Spearman captures monotonic ranking consistency.

\subsection{Experiment 2: Granularity, Spurious-Cue, and Human-Alignment Analysis}
\textbf{Main purpose:} Test whether personality evidence is local or distributed, whether model behavior depends on non-causal shortcuts (topic/slang cues), and whether model-identified important text aligns with human judgment.

\textbf{Procedure:}
\begin{itemize}
    \item Remove top-$k$ important words and re-evaluate.
    \item Remove top-$k$ important sentences and re-evaluate.
    \item Remove predefined topic-heavy words and stylistic markers, then re-evaluate.
    \item Compare with random-removal baselines matched for deletion size.
    \item For a sampled subset, collect human-annotated important words/sentences and compare these with model explanation outputs.
\end{itemize}

\textbf{Evaluation metrics:}
\begin{itemize}
    \item \textbf{Macro-F1 and accuracy drop} after each removal setting (tests information concentration).
    \item \textbf{Prediction flip rate} after spurious-cue removal (tests shortcut dependence).
    \item \textbf{Relative degradation gap} between sentence-removal and word-removal (tests local vs global encoding).
    \item \textbf{Human-model agreement} using top-$k$ overlap percentage and Jaccard similarity between human-selected and model-selected words/sentences (tests explanation alignment with human judgment).
\end{itemize}
These metrics directly map to RQ2, RQ3, and RQ4 and are reproducible via fixed seeds and fixed deletion budgets.

\clearpage



%%
%% The next two lines define the bibliography style to be used, and
%% the bibliography file.
\bibliographystyle{ACM-Reference-Format}
\bibliography{sample-base}
\appendix


\end{document}
\endinput
%%
%% End of file `sample-sigplan.tex'.
